\section{The model and the exact law}\label{sec:multistate-model}

We write $P^\mu_N$ in three forms. The \emph{run form} counts the switches,
and it gives all our exact statements. The \emph{refresh form} uses the
transition matrix $\sigma I+rJ$. When $p\ge1/d$ the chain keeps its state with
probability $\sigma$ and otherwise draws a fresh state uniformly from all $d$
states, possibly the same one. The \emph{ratio form} is the refresh form
divided by a vertex (so it measures a bin in units of a corner). For $p>1/d$ it is a sum of positive terms, and all our
asymptotics start from it.

For a positive count vector $n=(n_1,\ldots,n_k)$ with total $N$ and a state
$a$, let $W^{(a)}_n(j)$ be the number of words with counts $n$ and $j$ changes
that start with $a$, and let $W_n(j)=\sum_aW^{(a)}_n(j)$. Put
\[
S^{(a)}_n(u)=\sum_jW^{(a)}_n(j)u^j,\qquad
S_n(u)=\sum_aS^{(a)}_n(u)=\sum_jW_n(j)u^j .
\]
These polynomials do not depend on $d$. For a bin $n$ with some zero
coordinates, $S_n$ and $S^{(a)}_n$ mean the polynomials of its positive
coordinates.

\begin{proposition}[the generating function]\label{prop:dgf}
Put $y_i=x_it/(1-\sigma x_it)$ and $P^\mu_0(0,\ldots,0)=1$. As formal power
series,
\[
\sum_{N\ge0}\sum_{n_1+\cdots+n_d=N}P^\mu_N(n)\,x_1^{n_1}\cdots x_d^{n_d}\,t^N
=1+\frac{\sum_i\mu_iy_i}{1-r\sum_iy_i}.
\]
\end{proposition}

\begin{proof}
Let $F_i$ be the generating function of the words of length at least $1$ that
end in state $i$, each weighted by its probability, and let $S=\sum_iF_i$.
Conditioning on the last letter gives
$F_i=x_it\bigl(\mu_i+pF_i+r(S-F_i)\bigr)=x_it(\mu_i+\sigma F_i+rS)$, that is,
$F_i=y_i(\mu_i+rS)$. Summing over $i$ gives
$S=\sum_i\mu_iy_i+rS\sum_iy_i$, and the generating function is $1+S$.
\end{proof}

\begin{proposition}[the exact law]\label{prop:dpmf}
Let $N\ge1$, let $n$ be a bin, and put $F=\supp n$. Sums over $m$ run over the
integer vectors with $m_i=0$ for $i\notin F$ and $1\le m_i\le n_i$ for
$i\in F$, and $M=\sum_im_i$.
\begin{enumerate}[(a)]
\item (refresh form) With $0^0=1$,
\[
P^\mu_N(n)=\sum_m\Bigl(\sum_{i\in F}\mu_im_i\Bigr)\frac{(M-1)!}{\prod_im_i!}\,
r^{M-1}\sigma^{N-M}\prod_{i\in F}\binom{n_i-1}{m_i-1}.
\]
For the uniform start this reads
\[
P_N(n)=\frac1d\sum_m\frac{M!}{\prod_im_i!}\,
r^{M-1}\sigma^{N-M}\prod_{i\in F}\binom{n_i-1}{m_i-1}.
\]
\item (run form) If $p>0$, then $P^\mu_N(n)=p^{N-1}\sum_{a\in F}\mu_aS^{(a)}_n(u)$.
In particular $P_N(n)/V_N=S_n(u)$.
\item (ratio form) If $1/d<p<1$, so that $0<u<1$ and $v>0$, then
\begin{equation}\label{eq:simplex-positive-sum}
S_n(u)=(1-u)^{N-1}\sum_mM!\,v^{M-1}\prod_{i\in F}\frac1{m_i!}\binom{n_i-1}{m_i-1},
\end{equation}
which is a sum of positive terms.
\end{enumerate}
\end{proposition}

When $\sigma\ge0$, the refresh form has a direct reading. Cut the word into $M$
blocks at the $M-1$ refreshes. Each refresh lands on a prescribed state with
probability $r$, and each of the other $N-M$ steps keeps the state with
probability $\sigma$. If $m_i$ blocks carry state $i$, there are
$m_i(M-1)!/\prod_jm_j!$ orders of the block labels that start with $i$, and
$\binom{n_i-1}{m_i-1}$ ways to choose the lengths of the $m_i$ blocks of state
$i$. The proof below does not need $\sigma\ge0$.

\begin{proof}
(a) By the proof of Proposition~\ref{prop:dgf},
$S=\sum_{M\ge1}r^{M-1}\bigl(\sum_i\mu_iy_i\bigr)\bigl(\sum_iy_i\bigr)^{M-1}$. By the
multinomial theorem the coefficient of $\prod_iy_i^{m_i}$ in
$(\sum_i\mu_iy_i)(\sum_iy_i)^{M-1}$ is $\sum_i\mu_im_i(M-1)!/\prod_jm_j!$. With
$w_i=x_it$ we have $y_i=w_i/(1-\sigma w_i)$, and
$[w^j](w/(1-\sigma w))^m=\binom{j-1}{m-1}\sigma^{j-m}$ for $j\ge m\ge1$. For
$m=0$ only $j=0$ contributes. Multiplying the coefficients proves (a), and for
the uniform start $\sum_i\mu_im_i=M/d$.

(b) A word that starts with $a$ and has $j$ changes has probability
$\mu_ap^{N-1-j}r^j=\mu_ap^{N-1}u^j$. Sum over the words with counts $n$.

(c) Divide the uniform case of (a) by $V_N=p^{N-1}/d$. Since $r/p=u$ and
$\sigma/p=1-u$, we have
$r^{M-1}\sigma^{N-M}/p^{N-1}=u^{M-1}(1-u)^{N-M}=(1-u)^{N-1}v^{M-1}$.
\end{proof}

At $p=1/d$ we have $\sigma=0$ and only $M=N$ survives, which gives the
multinomial law for the uniform start. At $p=1$ only $M=1$ survives, which gives
mass $\mu_i$ at $Ne_i$. For $\sigma<0$ the terms in (a) can cancel, so for
numerical work we then use exact arithmetic or dynamic programming.

\subsection{Bessel bounds}

For $x\ge0$ the modified Bessel function of order $1$ is
$I_1(x)=\sum_{j\ge0}(x/2)^{2j+1}/(j!(j+1)!)$. Since
$\binom{2j+1}j\le4^j$, comparing coefficients gives $I_1(x)\le\sinh x\le e^x$.
We also use the expansion~\cite[Eqs.~10.40.1 and~10.17.1]{dlmf}
\[
I_1(x)=\frac{e^x}{\sqrt{2\pi x}}\Bigl(1-\frac3{8x}+O(x^{-2})\Bigr),
\qquad x\to\infty,
\]
and the function
\[
\Phi(x)=\sum_{j\ge0}\frac{x^j}{j!(j+1)!}=\frac{I_1(2\sqrt x)}{\sqrt x},
\qquad\Phi(0)=1 .
\]
Since $\binom{2j}j\le4^j$, we have $1/(j!(j+1)!)\le4^j/(2j)!$, and comparing term
by term with the even part of the series of $e^{2\sqrt x}$ gives
$\Phi(x)\le e^{2\sqrt x}$. Lemmas~\ref{lem:bessel-law}, \ref{lem:clipped-product}
and~\ref{lem:gaussian-phi-average} in Appendix~\ref{app:thresholds} give the
other facts about $\Phi$ that we use.
